% RESEARCH_DRAFT_SHA256: 83c3f42bef7b58e0646e562979af62e506240e55da2d80084b4aaf0edb3ba15b
% Layout and source-reference wrappers expanded for the web renderer.
\documentclass[11pt]{article}
\usepackage[margin=25mm]{geometry}
\usepackage{amsthm}
\usepackage{fontspec,unicode-math}
\setmainfont{Latin Modern Roman}
\setmathfont{Latin Modern Math}
\usepackage{amsmath,amssymb,mathtools,xurl,hyperref}
\hypersetup{colorlinks=true,urlcolor=blue,linkcolor=blue}
\setcounter{secnumdepth}{0}
\setlength{\parindent}{0pt}
\setlength{\parskip}{5pt}
\setlength{\emergencystretch}{3em}
\newcommand{\sref}[2]{\hyperlink{source:#1:#2}{[S#2]}}
\newcommand{\cataloguescope}[1]{\paragraph{Source record.}#1}
\newcommand{\floor}[1]{\left\lfloor#1\right\rfloor}
\newtheorem{proposition}{Proposition}
\newtheorem{lemma}[proposition]{Lemma}
\begin{document}
\paragraph{Research manuscript.} OpenAI \texttt{gpt-6-astra}, reasoning \texttt{ultra}; 11 October 2026.
\paragraph{Preserved source body.} The complete catalogue mathematical text and sources follow with layout and source links expanded, including source qualifications and its historical preparation notes. New research begins at the section headed Research attempt.

\section{Certified portfolios for a responsive Macaulay2 top level}\label{statement}

\subsection{Definitions and mathematical statement}

\paragraph{Shared notation.}Standard number systems are $\mathbb N,\mathbb Z,\mathbb Q,\mathbb R,\mathbb C$. Specialized notation and index ranges follow the cited source and the definitions below.


Concurrent user tasks let more than one calculation make progress in a session. A portfolio runs multiple algorithms for the same input and selects a valid result from one that completes; cancellation, resource sharing, and verification affect its behavior. Running calculations in the same process raises questions about shared state and safe interruption.
\paragraph{Record scope and qualifications.}
The catalogue title is retained as metadata. Named certificates, contracts, bounds or frameworks in that title are not asserted as results of the original source extract.
No task scheduler, correctness contract, resource budget, or speedup guarantee is specified. The workshop report records parallelization work; the prompt is historical.

Catalogue screening limitation: User-interface concurrency feature request rather than a mathematical target.

The source snapshot labels this record partially solved; this is catalogue metadata rather than independent verification of a mathematical result or the current completion of a project.
\paragraph{Statement recorded in the supplied source.}
In addition, the "top level" user interface would benefit from parellization.

What features at the user level should be parallelized? For which computations do we want to run multiple algorithms simultaneously and take whichever finishes first? Can we do other calculations in the same instance while one calculation is running? \href{http://aimpl.org/macaulay2efie/1/}{[S2]}
\paragraph{Source status.}
The supplied catalogue labels this source record partially solved; that metadata was checked against the supplied snapshot on 2026-10-11. The label is preserved for provenance and is not an independent verification of present open status or mathematical correctness.
\subsection{Short English statement}
Choose user-level Macaulay2 computations and workflows that should run concurrently.
\subsection{Sources}
\begin{itemize}
\item[S1] Ulam AI and the UnsolvedMath Contributors, UnsolvedMath: A Curated Collection of Open Mathematics Problems, catalogue entry. Catalogue curation and metadata: CC BY 4.0. The source extract is retained; definitions, English summaries and qualifications are editorial additions.
\url{https://huggingface.co/datasets/ulamai/UnsolvedMath}.
\item[S2] AIM Workshop Problem Lists, source linked by the catalogue.
\url{http://aimpl.org/macaulay2efie/1/}.
\item[S3] AIM workshop summary, Macaulay2: expanded functionality and improved efficiency (reports parallelization and sheaf-map implementation progress).
\url{https://aimath.org/pastworkshops/macaulay2efierep.pdf}.
\end{itemize}
\paragraph{Preparation and review.}
The original source and UnsolvedMath contributors retain attribution for the source extract. Codex (GPT-6; exact serving revision not exposed) assisted with definitions, English summaries, scope notes and formatting. No new human mathematical authorship or independently verified solution is asserted. Catalogue identity checked against the complete supplied registry and current repository. Distinct upstream records, including repeated questions and broad or damaged extracts, are retained for complete coverage of the source CS category. A repeated record is not a new independent conjecture. Complete mathematical-equivalence and current-literature review has not been performed. Preparation and final review follow contribution prompts C and D, with source gaps disclosed. Structural and rendering checks do not establish a new theorem or open status.
\label{end:statement}

\section{Research attempt}
\subsection{Literature, source context and scope}
The complete catalogue body above retains the original agenda, source status, and historical preparation. The AIMPL problem page could not be retrieved. The supplied complete source transcription was read, and the accessible primary AIM workshop announcement was checked for context. The announcement is not a substitute for verification of the inaccessible problem page. The catalogue's earlier proposed frameworks and research summaries are context, not established solutions credited to this attempt. The following is an explicitly limited investigation of one route; it neither substitutes a formalized toy model for the entire source agenda nor claims empirical or deployed-software results. Elementary auxiliary facts are not new original-target progress.

\subsection{A certified portfolio schedule under an explicit machine model}
The top-level question includes running several methods and accepting whichever finishes first. Consider $k\ge1$ deterministic algorithms on immutable copies of the same input. They can be paused after each primitive step. Their step counts include a sound verifier of the final answer to the same specified task. A rejected candidate does not count as completion. For this abstract model assume zero switching overhead and unit-cost steps; actual memory and scheduling costs are outside the bound.
\begin{proposition}
Cycling through the $k$ algorithms, one step each, obtains a verified answer within $kT$ total steps if at least one algorithm obtains one in $T$ steps when run alone.
\end{proposition}
\begin{proof}
After $r$ complete cycles, every still-needed algorithm has been given $r$ steps of its own execution, because immutable inputs and isolated states make interleaving preserve its sequential trajectory. By the end of cycle $T$, the successful algorithm has completed its $T$ steps and verification. At most $kT$ steps have then been executed; the scheduler may have stopped earlier. Acceptance is sound by the stipulated verifier for the common task.
\end{proof}
For runtimes $(3,100,7)$ and order 1,2,3, the first algorithm finishes on total step 7, below the bound 9. Isolation is material. If two tasks instead mutate one shared integer initially zero, with one operation incrementing and another doubling, the order increment-then-double gives 2 while double-then-increment gives 1. They cannot be treated as unchanged isolated sequential runs.
\subsection{Verification and remaining gaps}
The round-robin schedule and both shared-state orders were simulated. This is the standard dovetailing bound, not free speedup on one processor: it incurs a factor $k$ in this model. No responsive Macaulay2 top-level, cancellation protocol, sound task-specific verifier, memory bound or interactive shared-state policy was implemented. The original software feature decisions remain open here.

\section{References and source attribution}
Primary source: \url{http://aimpl.org/macaulay2efie/1/}. Attempted primary AIMPL page retrieval; failed with timeout or gateway error. The supplied full source transcription was read; the inaccessible live page was not independently inspected.

\url{https://aimath.org/pastworkshops/macaulay2efie.html}. Read the primary September 2023 workshop announcement.

\url{https://aimath.org/pastworkshops/macaulay2efierep.pdf}. Read the complete three-page primary workshop report. Existing threads/tasks, MPI work and parallelApply are acknowledged; no absence of all existing parallelism is inferred.

All retrievals were checked on 11 October 2026. Original contributor and catalogue attribution remains in the preserved Sources section. The elementary arguments below are self-contained; their exposition is not a claim of first discovery.


\paragraph{Integration review.} Codex (GPT-6; exact serving revision not exposed) reviewed the written arguments and expanded layout and source-reference wrappers. Repository metadata and the production web renderer were checked using Python 3.13 and JavaScript with MathJax 3.2.2. This model-assisted review is not independent human proof verification; no PDF compilation or formal proof checking is asserted.

\section{Completion Estimate}
\noindent\textbf{COMPLETION: 0\%}
No substantiated new progress toward the entire original target is claimed. The worked results and failed route are delimited above; a known result, illustrative model, or auxiliary obstruction is not counted as original-target completion.

\section{Final status and attribution}
\textbf{Attempt status: unresolved.} This describes this research attempt, not a new ruling on the source catalogue's status or on external literature claims.
The mathematics and this self-audit were generated by OpenAI \texttt{gpt-6-astra}, reasoning \texttt{ultra}, on 11 October 2026. Codex (GPT-6; exact serving revision not exposed) orchestrated the batch and contributed formatting/repository review. No human mathematical authorship, independent human proof checking, or formal verification is claimed. The historical preparation paragraphs in the reproduced catalogue body concern that body only.

Contribution prompt A was applied to the complete definition and supplied source record, followed by the adversarial checks of prompt B. Local reproduction files are \texttt{batch013/checks.py}, \texttt{computations.json}, and the per-ID audit JSON. They were not public links at preparation time. Numerical checks are finite and do not prove unrestricted results. No LaTeX engine was found locally; compilation was not performed by this research agent.

\end{document}
